\section{应用}
只要算符对 Lorentz 指标对称化且扣除全部迹，匹配的简单性即得到保证。
对任意 $n$ 写出无迹算符并不困难。
例如 $n=4$ 时
\be
\widehat{O}_{4444} = O_{4444} - \frac{3}{4}O_{\left\{\alpha \alpha 4 4 \right\}}
+ \frac{1}{16} O_{\left\{\alpha \alpha \beta \beta \right\}}\,.
\ee
$\widehat{O}_{4444}(t)$ 的强子矩阵元无需外部空间动量即可计算，
与重正化标度 $\mu$ 处 $\MS$ 方案的匹配由式~\eqref{eq:x_MS} 得到。

$\widehat{\rO}_n$ 约化强子矩阵元的连续极限记为
$A_n^h(t) = \left\langle x^{n-1} \right\rangle^h(t)$，
原则上需要从格点 QCD 模拟确定 $Z_\chi$。
但这可以避免：把二阶矩
$A_2^{h,\MS}(\mu) = \left\langle x \right\rangle_{\MS}^h(\mu)$
作为已经用任选方法算出的可观测量输入，
则 $n>2$ 的所有其他矩均可由下式计算
\be
\left\langle x^{n-1} \right\rangle_{\MS}^h(\mu) =
\left\langle x^{n-2} \right\rangle_{\MS}^h(\mu)
\frac{c_{n-1}(t,\mu)}{c_n(t,\mu)} R_n^h(t)\,,
\label{eq:ratio_n2}
\ee
其中 $c_n(t,\mu)$ 由式~(\ref{eq:cn_1l}-\ref{eq:finite}) 给出。比值
\be
R_n^h(t) = \frac{\left\langle x^{n-1} \right\rangle^{h}(t)}{\left\langle x^{n-2}\right\rangle^{h}(t)}\,,
\qquad n>2\,,
\label{eq:Rn}
\ee
具有有限的连续极限，因而无需重正化带流场。
此外，一旦格点理论经过非微扰改进，
这些比值只表现出 O($a^2$) 的标度破坏。
若数值上方便，还可以在式~\eqref{eq:ratio_n2}
的矩矩阵元计算中使用不同强子，
因为匹配系数与我们计算矩阵元所用强子无关。

本文提出的方法也可推广用于研究分布振幅，
但由于存在对前向矩阵元贡献为零的算符，
研究短流时间展开时须加小心~\cite{Bali:2017ude,RQCD:2019osh}。

梯度流此前已被用于在标量场论中提出局域光滑化的算符乘积展开~\cite{Monahan:2015lha}，
并用于定义准部分子分布（quasi-PDF）~\cite{Monahan:2016bvm}，
其单圈匹配见文献~\cite{Monahan:2017hpu}。
梯度流在准 PDF 计算中更现实的应用参见文献~\cite{HadStruc:2022yaw}。
原则上确定 PDF 的矩应能精确重构 PDF；
本文概述的 PDF 矩的确定能否与准 PDF 计算相结合，
是一个值得探索的有趣问题。
